As a rule, theses must be submitted in bound form. Usually comb or adhesive
binding is used. In this case, a so-called binding correction must be made. The
following details are taken from from \cite{dtk04.1:kohm:ausgleich}.

Since a part of the inner margin is covered by the binding and the bent paper,
the type area must be made a little narrower. This means that the proportions
of the opened work are visually coherent again. This distance binding
correction~$K$ is determined by the thickness of the book block book block~$d$ and
the width of the binding~$b$.
%
\begin{equation}
K = \dfrac{1}{2} \cdot \max(d, b)
\end{equation}
%
The book block thickness can be calculated from the grammage~$G$ and the number
of sheets~$n$. Here, the grammage in the unit $\unit{g / m^2}$ into the size
equation. The result of the formula corresponds to the thickness in
millimetres.
%
\begin{equation}
d = \dfrac{n \cdot G}{1000}
\end{equation}
%
A book block of 100~pages (that is 50~sheets for double-sided printing)
therefore has a thickness of \unit[4]{mm} for paper with a thickness of $\unit[80]{g / m^2}$.
However, a usual comb binding has a width of \unit[6]{mm}, so that the
correction must be increased to this value.

The binding correction can be set via the package~\ctanpkg{geometry}.
Listing~\ref{lst:tpl:bcor} shows the corresponding configuration option.

\lstinputlisting[language={[LaTeX]{TeX}},style=block,escapechar=\!,caption={set binding correction},label={lst:tpl:bcor}]{code/tpl_bcor.tex}

Although the number of pages is not fixed until the work is completed, the
binding correction should be determined at the beginning by a rough
calculation. A change in the binding correction will lead to a different line
and page break pagination, so that in this case certain fine-tuning may have to
be made again.

No binding correction is necessary for ring bindings.
